\StartFrontChapter{Schlusswort}\label{sec:14-schlusswort}
\begin{contentbox}
Ich wünsche allen viel Erfolg bei der Prüfung! Die ZSF soll bei Theoriefragen und beim Umsetzen von Programmieraufgaben helfen: Code-Skelette zeigen übertragbare Grundmuster. Fachliche Inspiration: Nina Burren, Georg Niggli und Eduard Fell. Fehler oder Ergänzungen gerne als Issue melden.

\ZSFgapS
\noindent\textbf{Links \& Ressourcen:}\\
\textbf{Aktuelle Version \& Hliddal Template:}\\
\href{https://garcon-studios.ch/zsf}{garcon-studios.ch/zsf}
\end{contentbox}

\begin{contentbox}
\ZSFkeyword*{HINWEIS ZUM DRUCKEN} \\
Beim Drucken \ZSFkeyword*{2 Seiten pro Blattseite} wählen — so werden die 16 PDF-Seiten zu den erlaubten 8 Druckseiten (entspricht 4 A4-Blättern). Je zwei PDF-Seiten bilden eine Druckseite (\texttt{1a}+\texttt{1b} bis \texttt{8a}+\texttt{8b}).
\end{contentbox}
